\section{7変数の代数的族と一般点の最小周期2}\label{sec:algebraic-family}
ここでは形式的な持上げから離れ、全パラメータで定義された代数的族を構成する。
基底体は引き続き標数零とする。基底を
\[
 S=\operatorname{Spec}k[u,v,r,s,a,b,c]\simeq\mathbb A^7
\]
とし、位置$i$で$\varepsilon=(-1)^i$と置く。積は隣接する一次生成空間の積を表す。
\begin{equation}\label{eq:seven-family}
\begin{aligned}
 YW+WY&=0,& ZW+WZ&=0,& ZY+YZ&=W^2,\\
 XW+WX&=uY^2+vZ^2+\varepsilon a(YZ-\tfrac12W^2),\\
 XY+YX&=W^2+Z^2+rY^2+\varepsilon(aWY+bWZ),\\
 XZ+ZX&=sY^2+\varepsilon cWY.
\end{aligned}
\end{equation}
この二次$\ZZ$-代数を$\mathcal C$と書く。
点$p_0=(0,1,0,0,0,0,0)$では$\mathcal C_{p_0}=C_0$である。
$u=r=s=a=b=0,\ v=1$に制限すると、既存の2周期族$A(c)$を得る。

\begin{theorem}[全ファイバーで正則な代数的族]\label{thm:algebraic-seven}
式\eqref{eq:seven-family}の$\mathcal C$は$S$上成分ごとに自由であり、
各$\mathcal C_{i,i+n}$の基底は
\[
 \{W^{h}Y^{j}Z^{\ell}X^{m}:h+j+\ell+m=n\}
\]
である。任意の幾何学的ファイバーは行Hilbert級数$(1-t)^{-4}$の
標準14641型quadratic Koszul AS正則$\ZZ$-代数である。
全ファイバーが2周期的である。ここには一般性条件や分母の非零条件を課さない。
\end{theorem}
\begin{proof}
$W<Y<Z<X$とし、逆順の二次語を順序付き語へ書き換える。
補助重みを$(1,2,2,4)$とし、重みが同じなら辞書式順序で比較する。
各規則の$-x_ax_b$以外の項は左辺より真に重みが小さく、
$-x_ax_b$も辞書式順序を小さくする。
従って全パラメータで還元は停止する。
重なりは$ZYW,XYW,XZW,XZY$の四つで、始点の偶奇を含めて8個である。
全ての差は$k[u,v,r,s,a,b,c]$上で恒等的に零となる。
再現コード\path{computations/verify_algebraic_family.py}は、特殊化や次数打切りを行わず
8個を多項式として検査する。Diamond lemma\cite[Theorem 1.2]{Bergman}から主張の自由基底を得る。

この恒等式には次の構造的説明もある。
\[
 B=k\langle W,Y,Z\rangle/(YW+WY,ZW+WZ,ZY+YZ-W^2),\quad T=YZ-\tfrac12W^2
\]
と置く。$W^2,Y^2,Z^2$は中心的で、$T$は$W$と可換、$Y,Z$と反可換である。
通常の次数1の導分$d$を
\[
 d(W)=aT,\qquad d(Y)=aWY+bWZ,\qquad d(Z)=cWY
\]
で定められる。最初の二関係の像は零であり、第三関係については
\[
 \{d(Z),Y\}+\{Z,d(Y)\}=2aWT=\{d(W),W\}
\]
となる。ここで$\{x,y\}=xy+yx$である。
一方、$\sigma_B(W,Y,Z)=(-W,-Y,-Z)$に対し、
\[
 p_W=uY^2+vZ^2,\quad p_Y=W^2+Z^2+rY^2,\quad p_Z=sY^2
\]
は、$\delta_p(fh)=\delta_p(f)h+\sigma_B(f)\delta_p(h)$を満たす$\sigma_B$-導分を定める。
各$p_j$が中心的なので関係の像は交換子の和として零になる。
$g^2=1$、$gW=-Wg,gY=-Yg,gZ=-Zg$を追加した環に
\[
 \sigma(g)=-g,\quad\sigma(W,Y,Z)=(-W,-Y,-Z),\quad
 \delta(g)=0,\quad\delta(j)=p_j+g d(j)
\]
を定めると、$\delta$は$\sigma$-導分である。
Ore拡大$Xh=\sigma(h)X+\delta(h)$の二つの冪等元$(1\pm g)/2$に沿った
展開が式\eqref{eq:seven-family}である。これも全パラメータでの基底を説明する。

AS正則性をHilbert級数だけから結論しない。
二次Gr\"obner基底からKoszul性が従う\cite[Theorem 5.3]{Priddy}。
そのbar複体のフィルトレーションによる証明は、各経路の始点・終点を固定すれば$\ZZ$-代数にも適用できる。
補助重みの随伴次数代数は、全ての相異なる生成元が反可換する四変数の量子多項式代数である。
二次双対$E$には逆重みのフィルトレーションを入れる。
二次関係の直交補空間の初期空間を取ることで、その随伴次数代数は
\[
 k[\xi_W,\xi_Y,\xi_Z,\xi_X]/(\xi_W^2,\xi_Y^2,\xi_Z^2,\xi_X^2)
\]
の対応する$\ZZ$-版からの商となる。
Koszul性とHilbert級数の逆数公式により$\dim E_{i,i+n}=\binom4n$なので、この全射は同型である。
随伴次数代数の相補次数積は完全である。
任意の非零な先頭項に相補単項式を掛ければ最高次数で非零になるから、
元の$E$の相補次数積も非退化であり、同じ次元より完全である。
従って局所Frobenius性が全ファイバーで成り立つ。
左右のKoszul複体をこの完全対合で双対化すると、単純加群に対するExtは
次数4だけに一次元存在し、左右AS条件と14641型解消を得る。
Koszul双加群解消の長さも4であり、大域次元は4となる。
係数が偶奇だけに依存するので厳密2移動は同型である。
\end{proof}

\begin{theorem}[同型類の7次元性と三つの交代方向]\label{thm:seven-image}
この族のindexed $\ZZ$-代数同型類の像は7次元である。
正確には、稠密開集合上でパラメータの同型類への写像は一般に有限対一であり、
この意味で有効パラメータは7個ある。
$p_0$における接方向は偶4方向と交代3方向に分かれ、いずれも同型で消えない。
交代3方向の前節の座標は
\begin{equation}\label{eq:seven-odd-plane}
 (o_0,o_1,o_2,o_3,o_4,o_5)=(0,-b,-a,-b,0,c).
\end{equation}
これは$Q=0$の最大等方3平面であり、本族はその全方向を代数的に同時に積分している。
\end{theorem}
\begin{proof}
位置0の一つの関係空間だけを取り、
$\phi:S\to\Gr(6,k^4\otimes k^4)$と
$G=\PGL_4\times\PGL_4$の作用を考える。
$p_0$での関係座標における群作用の微分は$[N_0\ N_1]$であり、階数は30である。
式\eqref{eq:seven-family}を七つのパラメータで微分した$60\times7$行列を$J$とすると
\[
 \rk[N_0\ N_1]=30,\qquad \rk[N_0\ N_1\ J]=37.
\]
同梱コードはこれらを有理数上で厳密に検証する。
従って滑らかな37次元の源$G\times S$からGrassmannianへの軌道飽和写像は
$p_0$で階数37を持ち、その像の閉包は37次元である。
一般の$G$軌道は30次元なので、一般軌道を分離する有理商の次元は7である。
別の言い方では、一般のパラメータに対し、その関係空間と$G$-同値なパラメータは有限個である。
裸のindexed同型は必ず位置0の関係空間の$G$-同値を誘導するので、
同型による同一視はこれ以上の正次元ファイバーを生まない。
全域モジュライstackの存在や全周期構造の分類は、この次元評価に不要である。

偶方向について$\rk[D_1(1)\ J_{u,v,r,s}]=19=15+4$、
交代方向について$\rk[D_1(-1)\ J_{a,b,c}]=18=15+3$となる。
商基底へ移すと式\eqref{eq:seven-odd-plane}を得る。
代入すれば$Q=0$であり、階数6の非退化二次形式の等方空間として次元3は最大である。
積分可能性は単なる接空間の主張ではなく、前定理の全$S$上の族によって与えられる。
\end{proof}

\subsection{隣接する点対応の射影配置による非1周期性}
関係を双線型化し、始点の座標を$[w:y:z:x]\in\mathbb P^3$とする。
次の$6\times4$行列が標点の座標に掛かる。
\begin{equation}\label{eq:seven-point-matrix}
 M_\varepsilon(w,y,z,x)=
 \begin{pmatrix}
 y&w&0&0\\z&0&w&0\\-w&z&y&0\\
 x+\varepsilon aw/2&-uy&-vz-\varepsilon ay&w\\
 -w&x-ry-\varepsilon aw&-z-\varepsilon bw&y\\
 0&-sy-\varepsilon cw&x&z
 \end{pmatrix}.
\end{equation}
$E_\varepsilon\subset\mathbb P^3$をその$4\times4$小行列式の零点スキームとする。
これは位置の偶奇に対応する点対応の左射影である。
1周期同型が存在すれば$E_+$と$E_-$は射影線型同型でなければならない。

\begin{theorem}[一般のファイバーは最小周期2]\label{thm:seven-period-two}
$S$の非空なZariski開集合上で、$\mathcal C$の最小周期は2である。
従って7次元の同型類の像の一般点は、どのgraded代数にも由来しない。
これは特定のgraded表示からの変形枝だけを排除する主張ではない。
\end{theorem}
\begin{proof}
$e_X=[0:0:0:1]$以外の点を$e_X$から$\mathbb P^2_{w,y,z}$へ射影する。
行列\eqref{eq:seven-point-matrix}の最初の三行の行列式は$-w(w^2+2yz)$である。
従って射影像は直線$L=(w=0)$と非特異二次曲線$D=(w^2+2yz=0)$の上にある。

$L$上で$[w:y:z]=[0:1:\ell]$とすると、標点の前三座標は$[0:1:-\ell]$であり、
必要十分条件は
\begin{equation}\label{eq:seven-line}
 v\ell^2+\varepsilon a\ell-u=0.
\end{equation}
$uv(a^2+4uv)\ne0$なら相異なる二点を得る。
残る座標は$x+x'=r-\ell^2,\ x'-x=s/\ell$から一意に決まる。

$D\setminus L$上では$[w:y:z]=[1:t:-1/(2t)]$で$t\ne0$と書け、
標点は$[1:-t:1/(2t)]$へ射影される。
最後の三関係から$x,x'$を消去すると
\begin{equation}\label{eq:seven-sextic}
 F_\varepsilon(t)=8st^6+8\varepsilon ct^5+4rt^4+4\varepsilon at^3
                   -4t^2-2\varepsilon bt+1=0
\end{equation}
を得る。$s\ne0$かつ$F_+$が重根を持たなければ、二次曲線上に相異なる六点がある。
それぞれの点で$x,x'$は一意に決まる。
この八点は、対応する一変数方程式の単根性から$E_\varepsilon$の被約点である。
一方$e_X$では最後の三行を用いて標点の三座標を局所的に消去でき、
残る方程式は二次以上である。従ってその接空間は三次元である。
上の条件では支持が有限なので$e_X$は非被約点であり、唯一の非被約支持点となる。

さらに、射影された八点のうち六点を通る二次曲線が$D$だけであるという条件を、
両符号について課す。これは単根の有限\'etale被覆上で有限個の行列式の非零として書ける開条件である。
この条件下で$E_+\simeq E_-$という射影線型同型は$e_X$を固定し、
商平面の射影変換を誘導する。その変換は$D$を保ち、残る二点を結ぶ$L$も保つ。
$D$上の座標$t$では、従って$t\mapsto\rho t$または$t\mapsto\rho/t$である。

$F_-(\rho t)$の根が$F_+(t)$の根に一致する場合、定数項は両方1なので両多項式は等しい。
$t^2$係数から$\rho^2=1$、$a\ne0$なら$t^3$係数から$\rho=-1$である。
反転型の場合は$t^6F_-(\rho/t)=F_+(t)/(8s)$となる。
$t^4,t^2$係数を比較して
\[
 -4\rho^2=\frac{r}{2s},\qquad 4r\rho^4=-\frac1{2s}
\]
を得るので、$r^3+8s=0$が必要である。これを非零と仮定すれば反転型はない。
残る$t\mapsto-t$は平面上で$[w:y:z]\mapsto[w:-y:-z]$であり、$L$上では恒等である。
しかし$a\ne0$なら式\eqref{eq:seven-line}の正負の二点集合は異なる。
従ってこの変換も両配置を対応させず、$E_+$と$E_-$は射影線型非同型である。

以上の開条件が非空であることを、次の有理数点で確認する。
\begin{equation}\label{eq:seven-rational-point}
 (u,v,r,s,a,b,c)=
 (13/3,-13/6,-1/4,1/2,13/2,13/5,-13/5).
\end{equation}
$F_+$の根は$1/5,1/2,1,5,-1,-1/2$、直線上の正符号の根は$1,2$、負符号では$-1,-2$である。
$r^3+8s=255/64$、$a^2+4uv=169/36$となる。
射影平面上の八点から六点を選ぶ28通りについて、二次単項式の評価行列を有理数上で計算すると、
行列式が零になるのは二次曲線上の六点を取った場合だけである。
その場合の階数は5であり、二次曲線は一意である。両符号とも同じ検査に成功する。
具体的な点、行列、根の検査は同梱コードに実装し、証明書に記録する。
従って全条件を満たす非空開集合が存在し、その上では1周期性が排除される。
全ファイバーが2周期的なので、そこでの最小周期は2である。
\end{proof}

\begin{remark}[何を決定し、何をまだ決定していないか]
本節は、形式級数の存在ではなく全$\mathbb A^7$上の多項式関係、全閉ファイバーの正則性、
一般点の非1周期性、同型類の像の次元7を与える。
ただしこの像が全モジュライの既約成分であるとは主張しない。
Clifford点の全九つの偶方向を含む族でもなく、残る交代方向の代数的積分、
全成分の次元、全1周期構造の分類は未解決である。
特に本節の7次元を、前節の14次元候補と取り違えてはならない。
\end{remark}

